% GENERATED FILE. Edit canonical lessons, then run npm run generate:books.
% canonical-source-hash: fnv1a64:350b3c7a872696b8
% canonical-lessons: GU-C554-ras, GU-C554-dadiya, GU-C554-ghughro, GU-C554-jhumvu, GU-C554-nacavvu, GU-R554-review-548-554, GU-R554-review-final-554

\chapter{Dancing All Night}
\label{ch:gu-dancing-all-night}
\hlchaptermodality{554}

\begin{chapteropening}
\textbf{By the end of this chapter:} \emph{I can talk about dancing raas and garba through the night.}

Complete the last lesson of chapter 554: I can talk about dancing raas and garba through the night.
\end{chapteropening}

\section[\texorpdfstring{\gu{રાસ} (rās)}{રાસ (rās)}]{\gu{રાસ} (rās) — raas (a stick dance)}

\label{lesson:GU-C554-ras}



\begin{quote}
Before the new one: say the Gujarati for to chime, then the Gujarati for to be itching to dance.
\end{quote}

\subsection*{You'll want to know: \gu{રાસ}}
\textbf{\gu{રાસ}} — \emph{rās} — \textquotedblleft{}raas (a stick dance)\textquotedblright{}.

\begin{quote}
\gu{જો} \gu{આજે} \gu{રાત્રે} \gu{વરસાદ} \gu{નહીં} \gu{પડે} \gu{તો} \gu{આપણે} \gu{રાસ} \gu{રમીશું}.
\end{quote}

\begin{quote}
\emph{jo āje rātre varsād nahī̃ paḍe to āpaṇe rās ramīśũ.}
\end{quote}

\begin{quote}
If it doesn't rain tonight, we'll dance raas.
\end{quote}

A masculine noun.

\begin{cousinweb}
\textbf{\gu{રાસ}} (\emph{rās}) is a tatsama, taken from Sanskrit \emph{rāsa}, a country dance of cowherds (Monier-Williams). Gujarati \textquotedblleft{}plays\textquotedblright{} it: \textbf{\gu{રાસ} \gu{રમવો}}.
\end{cousinweb}

\begin{grammarlens}[title={What you've built}]
One more word for saying what happens if.
\end{grammarlens}

\subsection*{Your turn}
\begin{itemize}
\raggedright
  \item \emph{Say it:} \emph{rās}
  \item \emph{Say it:} \emph{rās}, once more
  \item \emph{Say it:} \emph{jo āje rātre varsād nahī̃ paḍe to āpaṇe rās ramīśũ.} — If it doesn't rain tonight, we'll dance raas.
  \item \emph{Recall:} say the Gujarati for to chime, then the Gujarati for to be itching to dance, then say \emph{rās} again
\end{itemize}

\begin{tcolorbox}[breakable,colback=teal!4,colframe=teal!35!black,title={Before you move on}]
What does \gu{રાસ} mean? (\textquotedblleft{}Raas (a stick dance)\textquotedblright{}.) Say it once more.
\end{tcolorbox}
\section[\texorpdfstring{\gu{દાંડિયા} (dā̃ḍiyā)}{દાંડિયા (dā̃ḍiyā)}]{\gu{દાંડિયા} (dā̃ḍiyā) — dandiya sticks}

\label{lesson:GU-C554-dadiya}



\begin{quote}
Before the new one: say the Gujarati for to be itching to dance, then the Gujarati for raas.
\end{quote}

\subsection*{You'll want to know: \gu{દાંડિયા}}
\textbf{\gu{દાંડિયા}} — \emph{dā̃ḍiyā} — \textquotedblleft{}dandiya sticks\textquotedblright{}.

\begin{quote}
\gu{જો} \gu{તારી} \gu{પાસે} \gu{દાંડિયા} \gu{નહીં} \gu{હોય} \gu{તો} \gu{હું} \gu{તને} \gu{મારા} \gu{આપીશ}.
\end{quote}

\begin{quote}
\emph{jo tārī pāse dā̃ḍiyā nahī̃ hoy to hũ tane mārā āpīś.}
\end{quote}

\begin{quote}
If you don't have dandiya sticks, I'll give you mine.
\end{quote}

A masculine noun, used in the plural.

\begin{cousinweb}
\textbf{\gu{દાંડિયા}} (\emph{dā̃ḍiyā}) is built on \emph{dā̃ḍo}, \textquotedblleft{}a pole, a stick\textquotedblright{}, which CDIAL lists under Sanskrit \emph{daṇḍá}, \textquotedblleft{}a stick\textquotedblright{} (CDIAL 6128); dancers strike them together in the raas.
\end{cousinweb}

\begin{grammarlens}[title={What you've built}]
One more word for saying what happens if.
\end{grammarlens}

\subsection*{Your turn}
\begin{itemize}
\raggedright
  \item \emph{Say it:} \emph{dā̃ḍiyā}
  \item \emph{Say it:} \emph{dā̃ḍiyā}, once more
  \item \emph{Say it:} \emph{jo tārī pāse dā̃ḍiyā nahī̃ hoy to hũ tane mārā āpīś.} — If you don't have dandiya sticks, I'll give you mine.
  \item \emph{Recall:} say the Gujarati for to be itching to dance, then the Gujarati for raas, then say \emph{dā̃ḍiyā} again
\end{itemize}

\begin{tcolorbox}[breakable,colback=teal!4,colframe=teal!35!black,title={Before you move on}]
What does \gu{દાંડિયા} mean? (\textquotedblleft{}Dandiya sticks\textquotedblright{}.) Say it once more.
\end{tcolorbox}
\section[\texorpdfstring{\gu{ઘૂઘરો} (ghūghro)}{ઘૂઘરો (ghūghro)}]{\gu{ઘૂઘરો} (ghūghro) — an ankle bell}

\label{lesson:GU-C554-ghughro}



\begin{quote}
Before the new one: say the Gujarati for raas, then the Gujarati for dandiya sticks.
\end{quote}

\subsection*{You'll want to know: \gu{ઘૂઘરો}}
\textbf{\gu{ઘૂઘરો}} — \emph{ghūghro} — \textquotedblleft{}an ankle bell\textquotedblright{}.

\begin{quote}
\gu{જો} \gu{તું} \gu{પગમાં} \gu{ઘૂઘરા} \gu{પહેરીશ} \gu{તો} \gu{નાચતી} \gu{વખતે} \gu{મીઠો} \gu{અવાજ} \gu{આવશે}.
\end{quote}

\begin{quote}
\emph{jo tũ pagmā̃ ghūghrā paherīś to nāctī vakhte mīṭho avāj āvśe.}
\end{quote}

\begin{quote}
If you wear ankle bells, they'll ring sweetly as you dance.
\end{quote}

A masculine noun.

\begin{cousinweb}
\textbf{\gu{ઘૂઘરો}} (\emph{ghūghro}) comes from \emph{*ghuṅghura}, a form CDIAL reconstructs, and CDIAL lists the Gujarati word, \textquotedblleft{}a jingling bell\textquotedblright{} (CDIAL 4477).
\end{cousinweb}

\begin{grammarlens}[title={What you've built}]
One more word for saying what happens if.
\end{grammarlens}

\subsection*{Your turn}
\begin{itemize}
\raggedright
  \item \emph{Say it:} \emph{ghūghro}
  \item \emph{Say it:} \emph{ghūghro}, once more
  \item \emph{Say it:} \emph{jo tũ pagmā̃ ghūghrā paherīś to nāctī vakhte mīṭho avāj āvśe.} — If you wear ankle bells, they'll ring sweetly as you dance.
  \item \emph{Recall:} say the Gujarati for raas, then the Gujarati for dandiya sticks, then say \emph{ghūghro} again
\end{itemize}

\begin{tcolorbox}[breakable,colback=teal!4,colframe=teal!35!black,title={Before you move on}]
What does \gu{ઘૂઘરો} mean? (\textquotedblleft{}An ankle bell\textquotedblright{}.) Say it once more.
\end{tcolorbox}
\section[\texorpdfstring{\gu{ઝૂમવું} (jhūmvũ)}{ઝૂમવું (jhūmvũ)}]{\gu{ઝૂમવું} (jhūmvũ) — to sway (to music)}

\label{lesson:GU-C554-jhumvu}



\begin{quote}
Before the new one: say the Gujarati for dandiya sticks, then the Gujarati for an ankle bell.
\end{quote}

\subsection*{You'll want to know: \gu{ઝૂમવું}}
\textbf{\gu{ઝૂમવું}} — \emph{jhūmvũ} — \textquotedblleft{}to sway (to music)\textquotedblright{}.

\begin{quote}
\gu{જો} \gu{સંગીત} \gu{સારું} \gu{હશે} \gu{તો} \gu{બધા} \gu{ઝૂમશે}.
\end{quote}

\begin{quote}
\emph{jo sangīt sārũ haśe to badhā jhūmśe.}
\end{quote}

\begin{quote}
If the music is good, everyone will sway to it.
\end{quote}

Past (he …): \gu{તે} \gu{ઝૂમ્યો}, \emph{te jhūmyo}.

\begin{cousinweb}
\textbf{\gu{ઝૂમવું}} (\emph{jhūmvũ}) is to sway with the music, the body and the head together. CDIAL lists Hindi \emph{jhūmnā}, \textquotedblleft{}to shake\textquotedblright{}, with a query, under a reconstructed \emph{*kṣumbhati}, \textquotedblleft{}shakes\textquotedblright{} (CDIAL 3726); it gives no Gujarati form, so its history is not settled.
\end{cousinweb}

\begin{grammarlens}[title={What you've built}]
One more word for saying what happens if.
\end{grammarlens}

\subsection*{Your turn}
\begin{itemize}
\raggedright
  \item \emph{Say it:} \emph{jhūmvũ}
  \item \emph{Say it:} \emph{jhūmvũ}, once more
  \item \emph{Say it:} \emph{jo sangīt sārũ haśe to badhā jhūmśe.} — If the music is good, everyone will sway to it.
  \item \emph{Recall:} say the Gujarati for dandiya sticks, then the Gujarati for an ankle bell, then say \emph{jhūmvũ} again
\end{itemize}

\begin{tcolorbox}[breakable,colback=teal!4,colframe=teal!35!black,title={Before you move on}]
What does \gu{ઝૂમવું} mean? (\textquotedblleft{}To sway (to music)\textquotedblright{}.) Say it once more.
\end{tcolorbox}
\section[\texorpdfstring{\gu{નચાવવું} (nacāvvũ)}{નચાવવું (nacāvvũ)}]{\gu{નચાવવું} (nacāvvũ) — to make (someone) dance}

\label{lesson:GU-C554-nacavvu}



\begin{quote}
Before the new one: say the Gujarati for an ankle bell, then the Gujarati for to sway.
\end{quote}

\subsection*{You'll want to know: \gu{નચાવવું}}
\textbf{\gu{નચાવવું}} — \emph{nacāvvũ} — \textquotedblleft{}to make (someone) dance\textquotedblright{}.

\begin{quote}
\gu{જો} \gu{તું} \gu{ઢોલ} \gu{વગાડીશ} \gu{તો} \gu{તું} \gu{બધાને} \gu{નચાવીશ}.
\end{quote}

\begin{quote}
\emph{jo tũ ḍhol vagāḍīś to tũ badhāne nacāvīś.}
\end{quote}

\begin{quote}
If you play the drum, you'll get everyone dancing.
\end{quote}

Past, with \gu{તેણે} (he or she …): \gu{તેણે} \gu{બધાને} \gu{નચાવ્યા}, \emph{teṇe badhāne nacāvyā}.

\begin{cousinweb}
\textbf{\gu{નચાવવું}} (\emph{nacāvvũ}) is built from \textbf{\gu{નાચવું}}, \textquotedblleft{}to dance\textquotedblright{}, with the causative \emph{-āv-} and a shortened vowel. CDIAL files \textbf{\gu{નાચવું}} under Sanskrit \emph{nṛtyati}, \textquotedblleft{}dances\textquotedblright{} (CDIAL 7583).
\end{cousinweb}

\begin{grammarlens}[title={What you've built}]
One more word for saying what happens if.
\end{grammarlens}

\subsection*{Your turn}
\begin{itemize}
\raggedright
  \item \emph{Say it:} \emph{nacāvvũ}
  \item \emph{Say it:} \emph{nacāvvũ}, once more
  \item \emph{Say it:} \emph{jo tũ ḍhol vagāḍīś to tũ badhāne nacāvīś.} — If you play the drum, you'll get everyone dancing.
  \item \emph{Recall:} say the Gujarati for an ankle bell, then the Gujarati for to sway, then say \emph{nacāvvũ} again
\end{itemize}

\begin{tcolorbox}[breakable,colback=teal!4,colframe=teal!35!black,title={Before you move on}]
What does \gu{નચાવવું} mean? (\textquotedblleft{}To make (someone) dance\textquotedblright{}.) Say it once more.
\end{tcolorbox}
\section[(dialogue)]{Looking back: from If We Win the Match to Dancing All Night}

\label{lesson:GU-R554-review-548-554}



\begin{quote}
Say the Gujarati for a team and a captain.
\end{quote}

\subsection*{Your turn}
Say these aloud:

\begin{itemize}
\raggedright
  \item a team, a captain, a match-up, a winner, to win it for (a team)
  \item football, a goal (in a game), a score, to catch a ball, a medal
  \item an arrow, a target, to aim at, to pierce (a target), a runner
  \item training, a sports coach, music practice, to make (someone) run, a trophy
  \item a stage, acting, a spectator, to make ring, to get going
  \item a kettledrum, a shehnai, to sing out, to chime, to be itching to dance
  \item raas, dandiya sticks, an ankle bell, to sway, to make (someone) dance
\end{itemize}

\begin{tcolorbox}[breakable,colback=teal!4,colframe=teal!35!black,title={Before you move on}]
A team? (\textbf{\gu{ટીમ}}, \emph{ṭīm}.) A captain? (\textbf{\gu{કપ્તાન}}, \emph{kaptān}.) A match-up? (\textbf{\gu{મુકાબલો}}, \emph{mukāblo}.) A winner? (\textbf{\gu{વિજેતા}}, \emph{vijetā}.) To win it for (a team)? (\textbf{\gu{જિતાડવું}}, \emph{jitāḍvũ}.) Football? (\textbf{\gu{ફૂટબોલ}}, \emph{phūṭbol}.) A goal (in a game)? (\textbf{\gu{ગોલ}}, \emph{gol}.) A score? (\textbf{\gu{સ્કોર}}, \emph{skor}.) To catch a ball? (\textbf{\gu{ઝીલવું}}, \emph{jhīlvũ}.) A medal? (\textbf{\gu{ચંદ્રક}}, \emph{candrak}.) An arrow? (\textbf{\gu{તીર}}, \emph{tīr}.) A target? (\textbf{\gu{નિશાન}}, \emph{niśān}.) To aim at? (\textbf{\gu{તાકવું}}, \emph{tākvũ}.) To pierce (a target)? (\textbf{\gu{ભેદવું}}, \emph{bhedvũ}.) A runner? (\textbf{\gu{દોડવીર}}, \emph{doḍvīr}.) Training? (\textbf{\gu{તાલીમ}}, \emph{tālīm}.) A sports coach? (\textbf{\gu{કોચ}}, \emph{koc}.) Music practice? (\textbf{\gu{રિયાઝ}}, \emph{riyājh}.) To make (someone) run? (\textbf{\gu{દોડાવવું}}, \emph{doḍāvvũ}.) A trophy? (\textbf{\gu{ટ્રોફી}}, \emph{ṭrophī}.) A stage? (\textbf{\gu{મંચ}}, \emph{manc}.) Acting? (\textbf{\gu{અભિનય}}, \emph{abhinay}.) A spectator? (\textbf{\gu{દર્શક}}, \emph{darśak}.) To make ring? (\textbf{\gu{ગજાવવું}}, \emph{gajāvvũ}.) To get going? (\textbf{\gu{જામવું}}, \emph{jāmvũ}.) A kettledrum? (\textbf{\gu{નગારું}}, \emph{nagārũ}.) A shehnai? (\textbf{\gu{શરણાઈ}}, \emph{śarṇāī}.) To sing out? (\textbf{\gu{લલકારવું}}, \emph{lalkārvũ}.) To chime? (\textbf{\gu{રણકવું}}, \emph{raṇakvũ}.) To be itching to dance? (\textbf{\gu{થનગનવું}}, \emph{thanganvũ}.) Raas? (\textbf{\gu{રાસ}}, \emph{rās}.) Dandiya sticks? (\textbf{\gu{દાંડિયા}}, \emph{dā̃ḍiyā}.) An ankle bell? (\textbf{\gu{ઘૂઘરો}}, \emph{ghūghro}.) To sway? (\textbf{\gu{ઝૂમવું}}, \emph{jhūmvũ}.) To make (someone) dance? (\textbf{\gu{નચાવવું}}, \emph{nacāvvũ}.)
\end{tcolorbox}
\section[(dialogue)]{Dancing All Night — once more}

\label{lesson:GU-R554-review-final-554}



\begin{quote}
Say the Gujarati for raas and dandiya sticks.
\end{quote}

\subsection*{Your turn}
\begin{itemize}
\raggedright
  \item \emph{Say it:} raas, dandiya sticks, an ankle bell, to sway, to make (someone) dance
\end{itemize}

\begin{tcolorbox}[breakable,colback=teal!4,colframe=teal!35!black,title={Before you move on}]
Raas? (\textbf{\gu{રાસ}}, \emph{rās}.) Dandiya sticks? (\textbf{\gu{દાંડિયા}}, \emph{dā̃ḍiyā}.) An ankle bell? (\textbf{\gu{ઘૂઘરો}}, \emph{ghūghro}.) To sway? (\textbf{\gu{ઝૂમવું}}, \emph{jhūmvũ}.) To make (someone) dance? (\textbf{\gu{નચાવવું}}, \emph{nacāvvũ}.)
\end{tcolorbox}
